% ============================================================================
% SECTION 2 — RELATED WORK
% Paper budget: ~1.5 pages (a thesis chapter surveys; a paper positions).
% ReProAgent (arXiv:2607.09123), our format model, splits this as
% "2 Background and Motivation" with one subsection per neighbouring area and
% a closing paragraph that says what each area leaves open. Same shape here.
% Rule for every paragraph here: name the work, name what it established,
% name what it leaves open for us. If a paper does not shape our design or
% our claim, it belongs in the citation list, not in a sentence of its own.
% Sources to adapt: thesisproposal/main.tex §2 and the related-work table in
% webapp/research_data.py (its vs_tdbf column is the positioning argument).
% ============================================================================
\section{Related Work}
\label{sec:related}

\subsection{Bug Reproduction and Test Generation}
\label{sec:related-brt}
% TODO (~2 paragraphs). Adapt proposal §2.1: AthenaTest~\cite{Tufano2020AthenaTest},
% Qureshi'25~\cite{Qureshi2025TestCase}, LIBRO~\cite{Kang2023LIBRO} (few-shot,
% non-agentic), Cheng'25 BRT Agent~\cite{Cheng2025Agentic} (28% plausible BRTs,
% +30% fixes when fed to APR). Newer agentic systems: ReProAgent~\cite{ReProAgent2026}
% (58/70% on SWT-bench), SWE-Doctor~\cite{SWEDoctor2026} (multi-faceted BRTs as
% runtime diagnosis). Close on the gap we use: these generate tests as a
% deliverable; none makes the test a precondition for the repair stage.

\subsection{Agentic Automated Program Repair}
\label{sec:related-apr}
% TODO (~2 paragraphs). Adapt proposal §2.2–2.3: SWE-agent~\cite{Yang2024SWEagent}
% as the canonical reactive agent (~12%), MarsCode~\cite{Liu2024MarsCode},
% Fan et al.~\cite{Fan2023APRLLMs}, SWE-bench itself~\cite{Jimenez2023SWEbench}.
% Meng'24~\cite{Meng2024Empirical} is the converging motivation: localization and
% reproduction are THE bottlenecks. Add the APR traceability study
% \cite{Traceability2026} — its "shift-left" argument is ours in different words.

\subsection{Test-First Ordering in Agentic Repair}
\label{sec:related-ordering}
% TODO (~2 paragraphs, and the most important two in this section — this is the
% one work a reviewer will ask us about). Cheng'26 "Dynamic Cogeneration"
% \cite{Cheng2025Cogeneration}: 120 internal Google bugs, Gemini 2.5 Pro,
% TDD/TLD/Freeform selected by prompt instruction; Freeform won, TLD beat TDD on
% candidateBRT@k, and the agent defaults to test-last when left free.
% Our position, stated plainly and fairly:
%   (a) their TDD is a prompt nudge inside ONE agent; ours is a two-agent
%       architecture with a hard, executed validation gate;
%   (b) their top failure mode (the agent deletes the test as "temporary") is
%       impossible here by construction;
%   (c) private internal bugs vs. a public benchmark others can rerun;
%   (d) their failure mode #3 (fix overfits the BRT) is independently
%       reproduced by our verification-gap result (Section~\ref{sec:disc-verification}).
% Also TDFlow~\cite{Yuan2025TDFlow}: TDD framing, no controlled ordering experiment.
% Note Akita'26~\cite{Akita2026MethAI} is a practitioner account — frame it as such.

\subsection{Measuring Repair Outcomes}
\label{sec:related-eval}
% TODO (~1 short paragraph; merged from the thesis's separate evaluation
% section to keep this section to one column). Why we report observable
% outcomes (resolved-rate, held-out tests) instead of proxy similarity metrics:
% BLEU/CodeBLEU correlate poorly with correctness~\cite{Evtikhiev2023BLEU}.
% One clause on fault-localization tooling (GZoltar~\cite{Campos2012Gzoltar},
% iFixR~\cite{Koyuncu2019iFixR}) to mark what we deliberately do not use.

\subsection{Positioning}
\label{sec:related-gap}
% TODO (3–4 sentences, no more). The gap in one sentence: no controlled
% evaluation of test-first as a mandatory ARCHITECTURAL constraint against a
% matched reactive baseline on public data. Cheng'26 comes closest and differs
% on three axes — prompt-level enforcement, private dataset, cogeneration in one
% agent rather than a gated two-phase pipeline. End by pointing at
% Section~\ref{sec:approach}.
